%%%%%%%%%%%%%%%%%%%%%%%%%%%%%%%%%%%%%%%%%%%%%%%%%%%%%%%%%%%%%%%%%%%%%%%%
% Compel–ZipMix : ICLR 2025 submission draft with improved introduction
%%%%%%%%%%%%%%%%%%%%%%%%%%%%%%%%%%%%%%%%%%%%%%%%%%%%%%%%%%%%%%%%%%%%%%%%
% https://docs.google.com/presentation/d/1XoQ24_KofQUOeSxLE_lNVjYAw0lMtsm9spF6ertwL7s/edit?slide=id.g2dc59faafd3_0_323#slide=id.g2dc59faafd3_0_323

% https://chatgpt.com/g/g-p-67a3cf655c2081918746eaa9857f63a2-compel-pt-pre-training-via-compression-zip-fit/c/684f209f-3364-8001-97a9-8b1bf9b08397
\documentclass{article}
\usepackage{iclr2025_conference,times}
% Math commands for Compel-ZipMix paper

% Common sets
\newcommand{\R}{\mathbb{R}}
\newcommand{\N}{\mathbb{N}}
\newcommand{\Z}{\mathbb{Z}}

% Calligraphic letters
\newcommand{\cA}{\mathcal{A}}
\newcommand{\cB}{\mathcal{B}}
\newcommand{\cD}{\mathcal{D}}
\newcommand{\cF}{\mathcal{F}}
\newcommand{\cL}{\mathcal{L}}
\newcommand{\cM}{\mathcal{M}}
\newcommand{\cN}{\mathcal{N}}
\newcommand{\cO}{\mathcal{O}}
\newcommand{\cP}{\mathcal{P}}
\newcommand{\cS}{\mathcal{S}}
\newcommand{\cT}{\mathcal{T}}
\newcommand{\cV}{\mathcal{V}}
\newcommand{\cX}{\mathcal{X}}

% Operators
\DeclareMathOperator*{\argmin}{arg\,min}
\DeclareMathOperator*{\argmax}{arg\,max}
\DeclareMathOperator{\E}{\mathbb{E}}


\usepackage{hyperref}
\usepackage{url}
\usepackage{enumitem}
\usepackage{algorithm}
\usepackage{algpseudocode}
\usepackage{amsmath,amsfonts}

\title{Compel--ZipMix:\\
Compression-Aligned Mixtures for Efficient\\
Task-Aware Language-Model Pre-training}

\author{Brando Miranda$^{1}$\thanks{Equal contribution.}\;
        Elyas Obbad$^{1*}$\;
        Sanmi Koyejo$^{1}$\\
        $^{1}$Stanford University \\
        \texttt{\{bmiranda,eobbad,sanmi\}@stanford.edu}
}

%\iclrfinalcopy  % <-- uncomment for camera-ready
\begin{document}
\maketitle

%=========================  Abstract  ==============================%
\begin{abstract}
Mixture composition is a pivotal yet under-explored axis of
large-language-model pre-training.  
\textsc{DoReMi} learns domain weights via Group-DRO, but empirical
replications report accuracy losses: its coarse, hand-defined domains
blur signal and noise, and its worst-case objective begins from an
uninformative uniform prior—two assumptions at odds with the
No-Free-Lunch theorem \citep{wolpert1997no}.  
We present \emph{Compel-ZipMix}, which partitions the corpus into
compression-ratio buckets and assigns each bucket a weight derived from a
ZIP-based similarity to a small “AGI’’ validation suite.  
ZipMix includes a zero-cost \emph{Static} variant that samples directly
from this alignment prior and a \emph{DRO} variant that refines the prior
with a single Group-DRO pass.  
Using an X-parameter proxy on \textsc{The Pile},
ZipMix-Static outperforms a uniform mix by \textbf{Y\,pp}, while
ZipMix-DRO reaches the same score \textbf{X$\times$ faster};
both mixtures transfer unchanged to $\times10$ and $\times100$ models,
remaining below the three-order “dark-matter’’ scaling limit in the GPT-4
report.  
Because the prior depends solely on the chosen validation set, a single
CPU-only ZIP-FIT pass is enough to retarget the mixture toward any new
task suite, enabling \emph{practically free, domain-conditioned}
foundation checkpoints such as \textsc{MMLU+MiniF2F} or
\textsc{MMLU+FineWeb+ProofNet+MATH}.
\end{abstract}

%=========================  Introduction  ===========================%
\section{Introduction}\label{sec:intro}

The quality of a pre-training \emph{mixture} is as influential as token
count and parameter size.  
When high-entropy crawl dominates, large language models (LLMs) devote
compute to memorising noise, and downstream performance stalls.

\textsc{DoReMi} tackles mixture design with group distributionally robust
optimisation (Group-DRO).  
While its original study reported speed-ups, independent attempts often
observe that DoReMi funnels probability mass into noisy domains and can
reduce accuracy \citep{fan2023doge}.  
We trace this outcome to two linked assumptions:

\textbf{H\textsubscript{prior}} — \emph{No task prior.}  
DoReMi starts from a uniform weighting of coarse domains
(\emph{Wikipedia}, \emph{Books}, …).  
By the No-Free-Lunch theorem, a learner with no prior preference is free
to prefer entropy; the optimiser thus boosts whichever bucket maximises
loss—often noise.

\textbf{H\textsubscript{max}} — \emph{Over-pessimistic objective.}  
The inner maximisation always targets the highest excess-loss domain,
conflating “hard because useful’’ with “hard because irreducible’’ and
systematically over-weighting diversity.

\medskip
\textbf{Compel-ZipMix.}  
To test and remedy these hypotheses we introduce two steps:

1. \emph{Compression buckets} let the corpus self-partition by compression
   ratio, a model-free proxy for information density.

2. \emph{Alignment prior} assigns each bucket a weight
   $\alpha^{(0)}_m\propto\sum_{d\in B_m}s(d)+\tau$ with
   $s(d)=1-\text{NCD}_{\text{zip}}\bigl(d,\mathcal{V}\bigr)$ and
   $\tau=10^{-3}$, where $\mathcal{V}$ is a small validation suite
   (MMLU, GSM8K, MiniF2F, …).

\emph{ZipMix-Static} samples directly from $\boldsymbol{\alpha}^{(0)}$
(addressing H\textsubscript{prior});  
\emph{ZipMix-DRO} adds one Group-DRO refinement (tempering
H\textsubscript{max}).  
We cap model size jumps to $10^{3}$ to avoid the empirical
“dark-matter’’ deviation of scaling laws
\citep{gpt4report,schaeffer2025darkmatter}.

\medskip
\textbf{Empirical preview.}  
On \textsc{The Pile} and \textsc{FineWeb} a proxy of X parameters trained
with ZipMix-Static exceeds the uniform baseline by +Y\,pp on an AGI
suite; ZipMix-DRO adds Y\,pp and reaches baseline accuracy X× faster.
The mixture transfers unchanged to models $\times10$ and $\times100$
larger.

\medskip
\textbf{Free, task-conditioned checkpoints.}  
Because $\boldsymbol{\alpha}^{(0)}$ depends only on the validation corpus,
re-scoring the data with a new suite—say
\textsc{MMLU+MiniF2F} for maths or
\textsc{MMLU+FineWeb+ProofNet+MATH+Putnam-AXIOM} for formal reasoning—
requires a single ZIP-FIT sweep and an optional short proxy refresh,
giving any lab the ability to mint \emph{domain-specialised yet broadly
capable} foundation models on commodity hardware.

\medskip
\textbf{Contributions.}
\begin{itemize}[leftmargin=1.15em,topsep=0pt,itemsep=0pt]
\item We convert DoReMi’s rumoured failure into two explicit hypotheses
      and validate them empirically.
\item We introduce compression buckets plus alignment priors, yielding
      ZipMix-Static and ZipMix-DRO—cheap, robust, and task-aware.
\item We demonstrate X× faster convergence and +Y\,pp accuracy gains, and
      release code, bucket assignments, and scoring scripts for
      reproducibility.
\end{itemize}


%----------------------------------------------------------------------
\section{Method}\label{sec:method}

\subsection{Notation}\label{sec:notation}
Let $\mathcal{D}=\{d_1,\dots,d_N\}$ be the unlabeled corpus and
$\mathcal{V}$ a small validation suite whose tasks define the desired
prior.
After the compression–ratio split (Sec.~\ref{sec:buckets}) we obtain a
set of buckets
\[
  \mathcal{B} \;=\;\bigl\{B_1,\dots,B_{\lvert\mathcal{B}\rvert}\bigr\},
\]
where $\lvert\mathcal{B}\rvert$ denotes the number of buckets.
We write $\ell_{\theta}(x)$ for the per-token negative log-likelihood of a
language model with parameters $\theta$ and use
$\mathrm{NCD}_{\!\mathrm{zip}}(\cdot,\cdot)$ for the normalised
compression distance.

\subsection{Step 1: Compression buckets}\label{sec:buckets}
For each document $d\in\mathcal{D}$ compute the compression ratio
\begin{equation}
\mathrm{CR}(d)=
\frac{\text{bytes}_{\text{zip}}(d)}
     {\text{bytes}_{\text{raw}}(d)}.
\end{equation}
A sequence of thresholds
$0=c_0<c_1<\dots<c_{\lvert\mathcal{B}\rvert}=1$ induces the
buckets $B_m=\{d\!:\!\mathrm{CR}(d)\in[c_{m-1},c_m)\}$.

\subsection{Step 2: Alignment prior}\label{sec:prior}
Define the ZIP-FIT similarity
\[
  s(d)=1-\mathrm{NCD}_{\!\mathrm{zip}}\bigl(d,\mathcal{V}\bigr).
\]
The alignment weight for bucket $m$ is
\begin{equation}\label{eq:alpha0-b}
\alpha^{(0)}_m \;=\;
\frac{\sum_{d\in B_m}s(d)+\tau}
     {\sum_{B_j\in\mathcal{B}}\sum_{d\in B_j}s(d)
      +\lvert\mathcal{B}\rvert\,\tau},
\qquad
\tau=10^{-3},
\end{equation}
yielding $\boldsymbol{\alpha}^{(0)}\in\Delta^{\lvert\mathcal{B}\rvert-1}$.

\subsection{Idea 0: Alignment-Weighted Pre-Training (AWPT)}
Training minimises
\begin{equation}\label{eq:awpt-final}
\theta^\star\;=\;
\argmin_{\theta\in\Theta}\;
\sum_{B_m\in\mathcal{B}}
\alpha^{(0)}_m\;\frac{1}{\lvert B_m\rvert}
\sum_{x\in B_m} \ell_\theta(x).
\end{equation}
Stochastic optimisation uses hierarchical sampling:
sample $m\!\sim\!\text{Categorical}(\boldsymbol{\alpha}^{(0)})$, then
$x\!\sim\!\text{Uniform}(B_m)$.

\subsection{Idea 1: Loss-Reweighted AWPT}
Keep the dataloader uniform over $\mathcal{D}$ and scale each loss term by
$\alpha^{(0)}_{\,m(x)}$, where $m(x)$ maps a document to its bucket.

\subsection{Idea 2: ZipMix-DRO (One-Shot Refinement)}
Starting from $\boldsymbol{\alpha}^{(0)}$, run one Group-DRO pass on a
proxy model to obtain $\boldsymbol{\alpha}^{\star}$
(Alg.~\ref{alg:zipmixdro}) and then optimise
Eq.~\eqref{eq:awpt-final} with $\boldsymbol{\alpha}^{\star}$.

\subsection{Scaling constraint}\label{sec:scaling}
We limit the target model to at most
$10^{3}$\,\texttimes\ the proxy size, remaining below the empirical
“dark-matter’’ regime where scaling laws deviate
\citep{gpt4report,schaeffer2025darkmatter}.


%----------------------------------------------------------------------
\section{Experiments}\label{sec:experiments}
\textbf{Setup.}
Proxy size X; target sizes X$\,\times\,10$ and X$\,\times\,100$.
Datasets: \textsc{The Pile} and \textsc{FineWeb}.
Validation suite: MMLU-dev, GSM8K-dev, MiniF2F-dev.

\textbf{Main results.}
ZipMix-Static beats uniform by +X\,pp.
ZipMix-DRO gains an extra +Y\,pp and reaches baseline accuracy
X$\times$ sooner.

\textbf{Ablations.}
Without the alignment prior, optimisation collapses into the noisiest
bucket; replacing CR buckets with hand domains reproduces DoReMi’s failure.

%----------------------------------------------------------------------
\section{Dark-Matter Scaling}\label{sec:darkmatter}
Scaling laws fit a power curve up to $\approx10^{3}$ scale; beyond that the
loss deviates unpredictably, a phenomenon termed \emph{dark matter}
\citep{schaeffer2025darkmatter}.  All ZipMix evaluations remain below this
threshold.

%----------------------------------------------------------------------
\section{Discussion and Limitations}\label{sec:discussion}
Our method depends on the relevance of the chosen validation suite and on the
language neutrality of compression ratio.  Multilingual corpora and radically
different downstream tasks may require re-scoring.  Group-DRO adds
X\,\% extra proxy compute; ZipMix-Static offers a cheaper fallback.

\bibliography{zipmix_refs}
\bibliographystyle{iclr2025_conference}
\end{document}
